% ===========================================================================
\chapter{Matrizes de Rastreabilidade Estrutural e Arquitetural}
\label{cap_matrizes_rastreabilidade}
% ===========================================================================

A rastreabilidade estrutural estabelece a comprovação formal de que todas as classes, interfaces e serviços do SIGEA-GTT implementam de forma direta os requisitos de software e suportam integralmente as transações dos casos de uso, alinhando-se à persistência físico-relacional.

% ===========================================================================
\section{Matriz de Rastreabilidade: Classes vs. Requisitos Funcionais}
\label{sec_matriz_classes_rf}
% ===========================================================================

A Tabela \ref{tab_rastreabilidade_classes_rf} mapeia os vinte e cinco requisitos funcionais (RF01 a RF25) para as classes e serviços do backend que realizam a lógica de negócio:

\begin{table}[!ht]
\centering
\caption{Matriz de Rastreabilidade entre Classes de Software e Requisitos Funcionais}
\label{tab_rastreabilidade_classes_rf}
\footnotesize
\begin{tabularx}{\textwidth}{p{1.8cm} p{1.4cm} p{4.8cm} X}
\toprule
\textbf{Requisito} & \textbf{Módulo} & \textbf{Classes e Interfaces Envolvidas} & \textbf{Papel Estrutural na Implementação} \\
\midrule
\texttt{RF01} & M1 & \texttt{AuthService}, \texttt{UfsEmailValidator} & Validação de credenciais e domínio institucional. \\
\texttt{RF02} & M1 & \texttt{AuthService}, \texttt{User} & Forçamento de troca de senha no primeiro login. \\
\texttt{RF03} & M1 & \texttt{AuthService}, \texttt{PasswordEncoder} & Alteração voluntária de senha no perfil do usuário. \\
\texttt{RF04} & M1 & \texttt{UserService}, \texttt{UserRole}, \texttt{User} & Segregação dos perfis de acesso RBAC. \\
\texttt{RF05} & M1 & \texttt{UserService}, \texttt{UserCreateDTO} & Obrigatoriedade de matrícula apenas para estudantes. \\
\texttt{RF06} & M2 & \texttt{GttModuleService}, \texttt{GttModule} & CRUD e parametrização dos 6 módulos canônicos. \\
\texttt{RF07} & M2 & \texttt{GttTriggerService}, \texttt{GttTrigger} & Gestão taxonômica dos 53 gatilhos clínicos do IHI. \\
\texttt{RF08} & M2 & \texttt{GttTriggerController}, \texttt{GttDTO} & Catálogo didático de consulta rápida com busca. \\
\texttt{RF09} & M2 & \texttt{HarmSeverityService}, \texttt{HarmSeverity} & Mapeamento e classificação de gravidade NCC MERP. \\
\texttt{RF10} & M3 & \texttt{AcademicClassService}, \texttt{Class} & Criação de turmas com docência titular e formato. \\
\texttt{RF11} & M3 & \texttt{AcademicClassService}, \texttt{UserRepository} & Associação N:N de discentes na turma. \\
\texttt{RF12} & M3 & \texttt{ActivityService}, \texttt{Activity} & Alocação de atividade prática com prazo limite. \\
\texttt{RF13} & M3 & \texttt{AcademicClassService}, \texttt{SubmRepo} & Bloqueio de encerramento de turma com pendências. \\
\texttt{RF14} & M4 & \texttt{Activity}, \texttt{ClinicalCaseData} & Estrutura JSONB de prontuário simulado de alta fidelidade. \\
\texttt{RF15} & M4 & \texttt{AuditSessionService}, \texttt{Activity} & Cronômetro de 20 minutos com pausas e avisos. \\
\texttt{RF16} & M4 & \texttt{IdentifiedTriggerData}, \texttt{Submission} & Vinculação de gatilho com evidência e gravidade. \\
\texttt{RF17} & M5 & \texttt{IshikawaData}, \texttt{QualityToolsData} & Diagrama de causa e efeito nos 6 vetores causais. \\
\texttt{RF18} & M5 & \texttt{GutItemData}, \texttt{QualityToolsData} & Matriz GUT com cálculo de score de prioridade. \\
\texttt{RF19} & M5 & \texttt{FiveWTwoHItemData}, \texttt{QualityTools} & Plano de ação estruturado nas 7 dimensões 5W2H. \\
\texttt{RF20} & M5 & \texttt{PdcaData}, \texttt{QualityToolsData} & Ciclo PDCA para melhoria contínua hospitalar. \\
\texttt{RF21} & M4, M5 & \texttt{ActivitySubmission}\newline\texttt{Service}, \texttt{Subm} & Submissão atômica em transação ACID única. \\
\texttt{RF22} & M6 & \texttt{ActivitySubmission}\newline\texttt{Service}, \texttt{Subm} & Atribuição de nota de 0 a 10 e parecer pedagógico. \\
\texttt{RF23} & M6 & \texttt{ClassDashboardService}, \texttt{Subm} & Comparação com gabarito e cálculo de sensibilidade. \\
\texttt{RF24} & M6 & \texttt{ClassDashboardService}, \texttt{GttMetrics} & Consolidação de taxas por 1.000 pacientes-dia. \\
\texttt{RF25} & M6 & \texttt{ClassDashboard}\newline\texttt{Controller}, \texttt{CSV} & Exportação de relatórios em formatos estruturados. \\
\bottomrule
\end{tabularx}
\fonte{Elaborado pelo autor (2026).}
\end{table}

% ===========================================================================
\section{Matriz de Rastreabilidade: Classes vs. Casos de Uso}
\label{sec_matriz_classes_uc}
% ===========================================================================

A Tabela \ref{tab_rastreabilidade_classes_uc} correlaciona as classes executivas aos vinte e dois casos de uso (UC01 a UC22) do SIGEA-GTT:

\begin{table}[!ht]
\centering
\caption{Matriz de Rastreabilidade entre Classes e Casos de Uso}
\label{tab_rastreabilidade_classes_uc}
\footnotesize
\begin{tabularx}{\textwidth}{p{2.2cm} p{4.8cm} X}
\toprule
\textbf{Caso de Uso} & \textbf{Classes Controladoras e de Serviço} & \textbf{Entidades e DTOs Manipulados} \\
\midrule
\texttt{UC01}--\texttt{UC03} & \texttt{AuthController}, \texttt{AuthService} & \texttt{User}, \texttt{LoginRequestDTO}, \texttt{LoginResponseDTO} \\
\texttt{UC04} & \texttt{UserController}, \texttt{UserService} & \texttt{User}, \texttt{UserCreateDTO}, \texttt{UserResponseDTO} \\
\texttt{UC05}, \texttt{UC06} & \texttt{GttModuleController}, \texttt{GttTriggerController} & \texttt{GttModule}, \texttt{GttTrigger}, \texttt{HarmSeverity} \\
\texttt{UC07}--\texttt{UC10} & \texttt{AcademicClassController}, \texttt{ClassService} & \texttt{AcademicClass}, \texttt{User}, \texttt{ClassCreateDTO} \\
\texttt{UC11}--\texttt{UC13} & \texttt{ActivityController}, \texttt{SubmissionService} & \texttt{Activity}, \texttt{ClinicalCaseData}, \texttt{TriggerData} \\
\texttt{UC14}--\texttt{UC17} & \texttt{ActivitySubmission}\newline\texttt{Service}, \texttt{QualityData} & \texttt{IshikawaData}, \texttt{GutItemData}, \texttt{FiveWTwoHData} \\
\texttt{UC18} & \texttt{ActivitySubmission}\newline\texttt{Controller}, \texttt{SubmService} & \texttt{ActivitySubmission}, \texttt{SubmissionCreateDTO} \\
\texttt{UC19}, \texttt{UC20} & \texttt{ActivitySubmission}\newline\texttt{Controller}, \texttt{SubmService} & \texttt{ActivitySubmission}, \texttt{SubmissionGradeDTO} \\
\texttt{UC21}, \texttt{UC22} & \texttt{ClassDashboardController}, \texttt{DashboardService} & \texttt{GttMetricsDTO}, \texttt{ClassDashboardDTO} \\
\bottomrule
\end{tabularx}
\fonte{Elaborado pelo autor (2026).}
\end{table}

% ===========================================================================
\section{Matriz de Padrões de Projeto Aplicados}
\label{sec_matriz_padroes_projeto}
% ===========================================================================

A Tabela \ref{tab_padroes_projeto_classes} cataloga os padrões de projeto adotados na arquitetura de software orientada a objetos:

\begin{table}[!ht]
\centering
\caption{Catálogo de Padrões de Projeto Aplicados no SIGEA-GTT}
\label{tab_padroes_projeto_classes}
\footnotesize
\begin{tabularx}{\textwidth}{p{3.0cm} p{2.5cm} p{4.2cm} X}
\toprule
\textbf{Padrão de Projeto} & \textbf{Classificação} & \textbf{Classes Representativas} & \textbf{Motivação e Benefício Arquitetural} \\
\midrule
\textbf{Repository} & Corporativo & \texttt{UserRepository}, \texttt{ClassRepository} & Abstração de persistência e isolamento do PostgreSQL. \\
\textbf{Service Layer} & Corporativo & \texttt{AuthService}, \texttt{ClassService} & Encapsulamento de regras clínicas e transações ACID. \\
\textbf{Data Transfer Object} & Corporativo & \texttt{UserCreateDTO}, \texttt{LoginResponseDTO} & Desacoplamento entre entidade interna e payload JSON. \\
\textbf{Mapper} & Estrutural & \texttt{UserMapper}, \texttt{TriggerMapper} & Conversão estática e performática via MapStruct. \\
\textbf{Singleton} & Criacional & Spring Beans (\texttt{@Service}, \texttt{@Controller}) & Instâncias únicas compartilhadas sem estado interno. \\
\textbf{Dependency Injection} & Estrutural & Spring Framework (\texttt{@Autowired}) & Baixo acoplamento e suporte integral a testes unitários. \\
\textbf{Interceptor Filter} & Comportamental & \texttt{JwtAuthentication}\newline\texttt{Filter}, \texttt{AuthInterceptor} & Validação transversal de autenticação HTTP em pipeline. \\
\bottomrule
\end{tabularx}
\fonte{Elaborado pelo autor (2026).}
\end{table}
